% Part:sets-functions-relations
% Chapter: size-of-sets
% Section: comparing-sizes

\documentclass[../../../include/open-logic-section]{subfiles}

\begin{document}

\olfileid{sfr}{siz}{car}

\olsection{Verzamelingen met verschillende aantallen elementen en de stelling van Cantor}

\begin{explain}
We hebben precies gezegd wat het betekent dat twee verzamelingen evenveel
elementen hebben. We kunnen ook precies zeggen wanneer de ene verzameling
minder elementen heeft dan de andere. Voor onze definitie van ``kleiner dan
(of even groot als)'' gebruiken we geen !!a{bijection} tussen de
verzamelingen, maar !!a{injection} van de eerste verzameling naar de tweede.
Als zo'n functie bestaat, heeft de eerste verzameling hoogstens evenveel
elementen als de tweede. Bij !!a{injection} krijgen verschillende invoeren
verschillende beelden. De functie neemt daardoor minstens evenveel
verschillende waarden aan als er !!{element}s in het domein zitten. Geen twee
!!{element}s uit het domein worden namelijk op hetzelfde !!{element} afgebeeld.
\end{explain}

\begin{defn}
$A$ is \emph{niet groter dan}~$B$, geschreven als $\cardle{A}{B}$, precies
als er !!a{injection} $f \colon A \to B$ bestaat.
\end{defn}

Deze relatie is duidelijk reflexief en transitief, maar niet symmetrisch
(dit moet je zelf bewijzen). We kunnen nu ook vastleggen wanneer de ene
verzameling (strikt) kleiner is dan de andere. 

\begin{defn}
$A$ is \emph{kleiner dan}~$B$, geschreven als $\cardless{A}{B}$, precies als
er !!a{injection}~$f\colon A \to B$ bestaat, maar geen
!!{bijection}~$g\colon A \to B$. Er geldt dus $\cardle{A}{B}$ en
$\cardneq{A}{B}$.
\end{defn}

Ook deze relatie is duidelijk irreflexief en transitief. (Ook dit moet je
zelf bewijzen.) Met deze notatie kunnen we zeggen dat een verzameling $A$
!!{enumerable} is precies als $\cardle{A}{\Nat}$, en dat $A$
!!{nonenumerable} is precies als $\cardless{\Nat}{A}$. Daarmee kunnen we
\oliflabeldef{sfr:siz:nen-alt:thm:nonenum-pownat}{%
\olref[sfr][siz][nen-alt]{thm:nonenum-pownat}
ook zo formuleren:
$\cardless{\Nat}{\Pow{\Nat}}$}{%
\olref[sfr][siz][nen]{thm:nonenum-pownat}
ook zo formuleren: $\cardless{\PosInt}{\Pow{\PosInt}}$}. Sterker nog,
\citet{Cantor1892} bewees dat deze laatste uitspraak \emph{voor elke
verzameling} geldt:

\begin{thm}[Cantor]\ollabel{thm:cantor}
$\cardless{A}{\Pow{A}}$, voor elke verzameling $A$.
\end{thm}

\begin{proof}
De functie $f(x) = \{x\}$ is !!a{injection} $f \colon A \to \Pow{A}$.
Als $x \neq y$, bevatten de twee verzamelingen met maar één element niet
dezelfde elementen. Volgens extensionaliteit geldt dan ook
$\{x\} \neq \{y\}$, en dus $f(x) \neq f(y)$. Daarom geldt
$\cardle{A}{\Pow{A}}$. 

\begin{editorial}
We geven het langere bewijs als \olref[nen]{sec} in de tekst staat. Anders
geven we het kortere bewijs dat bij \olref[nen-alt]{sec} hoort.
\end{editorial}

\oliflabeldef{sfr:siz:nen:sec}{% 
  We laten nu zien dat er geen surjectieve functie~$g\colon A \to
  \Pow{A}$ kan bestaan, laat staan een bijectieve functie. Daaruit volgt dat
  $\cardneq{A}{\Pow{A}}$. Neem namelijk een functie
  $g\colon A \to \Pow{A}$. Omdat $g$ totaal is, krijgt elke $x \in A$ een
  deelverzameling $g(x) \subseteq A$ als waarde. We kunnen laten zien dat $g$
  niet surjectief kan zijn. Daarvoor maken we een deelverzameling
  $\overline{A} \subseteq A$ die volgens haar definitie niet voorkomt tussen
  de waarden die~$g$ werkelijk aanneemt. Neem
  \[
  \overline{A} = \Setabs{x \in A}{x \notin g(x)}.
  \]
  Omdat $g(x)$ voor elke $x \in A$ bestaat, legt deze regel duidelijk één
  deelverzameling $\overline{A}$ van~$A$ vast. Toch komt zij niet voor tussen
  de waarden die~$g$ aanneemt. Neem een willekeurige $x \in A$. We laten zien
  dat $\overline{A} \neq g(x)$. Als $x \in g(x)$, dan geldt
  $x \notin g(x)$ niet. Uit de definitie van~$\overline{A}$ volgt dan dat
  $x \notin \overline{A}$. Als $x \notin g(x)$, dan geldt ook $x \in A$,
  want zo is dit element gekozen. Daarmee voldoet het aan de regel waarmee
  $\overline{A}$ is bepaald, en dus geldt $x \in \overline{A}$. De twee
  gevallen laten voor de verzameling~$\overline{A}$ en elke $x \in A$ zien:
  $x \in g(x)$ precies als
  $x \notin \overline{A}$. Dus $g(x) \neq \overline{A}$. Anders gezegd:
  $\overline{A}$ komt niet voor tussen de waarden die~$g$ aanneemt. Dat
  spreekt de aanname tegen dat~$g$ surjectief is.}{We hoeven alleen nog te
  laten zien dat $\cardneq{A}{\Pow{A}}$. Neem, om een tegenspraak te krijgen,
  aan dat $\cardeq{A}{\Pow{A}}$. Er is dan !!{bijection}
  $g \colon A \to \Pow{A}$. Kijk nu naar:
  \[
    D = \Setabs{x \in A}{x \notin g(x)}
  \]
  Omdat $D \subseteq A$, geldt $D \in \Pow{A}$. Omdat $g$ !!a{bijection}
  is, is er een $y \in A$ waarvoor $g(y) = D$. Maar dan geldt:
  \[
    y \in g(y) \text{ precies als } y \in D \text{ precies als } y \notin g(y).
  \]
  Dat is een tegenspraak; dus $\cardneq{A}{\Pow{A}}$.}{}
\end{proof}

\begin{explain}
\oliflabeldef{sfr:siz:nen:thm:nonenum-pownat}{Het helpt om het bewijs van
  \olref{thm:cantor} te vergelijken met dat van
  \olref[nen]{thm:nonenum-pownat}. Daar bewezen we het volgende. Neem een
  willekeurige lijst $Z_1$, $Z_2$, \dots, van deelverzamelingen
  van~$\PosInt$. We kunnen dan een verzameling~$\overline{Z}$ van getallen
  maken die zeker niet op de lijst staat. Dat weten we omdat voor elke
  $n \in \PosInt$ geldt: $n \in Z_n$ precies als $n \notin \overline{Z}$.
  Er is daardoor altijd een getal dat !!a{element} is van een van $Z_n$ en
  $\overline{Z}$, maar geen !!{element} van de andere. We doen hier
  hetzelfde. Alleen betekent index~$n$ nu: een !!{element} van~$A$, en niet
  een getal uit~$\PosInt$. We hebben $\overline{A}$ zo bepaald dat deze
  verzameling verschilt van~$g(x)$ voor elke $x \in A$: $x \in g(x)$ precies als
  $x \notin \overline{A}$. Ook nu vinden we steeds !!a{element} van~$A$.
  Dat is !!a{element} van precies één van de twee verzamelingen: $g(x)$ of
  $\overline{A}$. $\overline{Z}$ kan daardoor niet op de lijst $Z_1$, $Z_2$, \dots,
  staan. Om dezelfde reden komt $\overline{A}$ niet voor tussen de waarden
  die~$g$ werkelijk aanneemt.}{}

\oliflabeldef{sfr:siz:nen-alt:thm:nonenum-pownat}{Het helpt om het bewijs van
  \olref{thm:cantor} te vergelijken met dat van
  \olref[nen-alt]{thm:nonenum-pownat}. Daar bewezen we het volgende. Neem een
  willekeurige lijst $N_0$, $N_1$, $N_2$, \dots, van deelverzamelingen
  van~$\Nat$. We kunnen dan een verzameling~$D$ van getallen maken die zeker
  niet op de lijst staat. Dat weten we omdat $n \in N_n$ precies als
  $n \notin D$, voor elke $n \in \Nat$. We doen hier hetzelfde. Alleen
  betekent index~$n$ nu: een !!{element} van~$A$, en niet een getal uit~$\Nat$.
  We hebben $D$ zo bepaald dat deze verzameling verschilt van~$g(x)$ voor elke $x \in A$:
  $x \in g(x)$ precies als $x \notin D$.}{}

Dit bewijs lijkt ook op het bewijs van de paradox van Russell,
\olref[sfr][set][rus]{thm:russells-paradox}. De stelling van Cantor bracht
Russell zelfs op het idee voor zijn eigen paradox.
\end{explain}

\begin{prob}
  Laat zien dat er geen !!a{injection} $g\colon \Pow{A} \to A$ kan bestaan,
  welke verzameling~$A$ we ook nemen. Tip: neem aan dat
  $g\colon \Pow{A} \to A$ !!{injective} is. Kijk naar
  $D = \Setabs{g(B)}{B \subseteq A \text{ en } g(B) \notin B}$. Stel
  $x = g(D)$. Gebruik het feit dat $g$ !!{injective} is om een tegenspraak te
  krijgen. 
\end{prob}
\end{document}
